\documentclass[11pt,a4paper]{article}
\usepackage[margin=25mm]{geometry}
\usepackage{lmodern}
\usepackage[T1]{fontenc}
\usepackage{booktabs}
\usepackage{array}
\usepackage{enumitem}
\usepackage{fancyhdr}
\usepackage{titlesec}
\usepackage[hidelinks]{hyperref}

\linespread{1.08}
\setlist{itemsep=2pt,parsep=0pt,topsep=3pt}
\titlespacing*{\section}{0pt}{14pt}{6pt}
\titlespacing*{\subsection}{0pt}{10pt}{4pt}
\titleformat{\section}{\large\bfseries}{\thesection.}{0.6em}{}
\titleformat{\subsection}{\bfseries}{}{0em}{}

\pagestyle{fancy}
\fancyhf{}
\lhead{\small\itshape Lesson-plan-driven generation of personalised study packs}
\rhead{\small 5 August 2026}
\cfoot{\small\thepage}
\renewcommand{\headrulewidth}{0.4pt}

\newcommand{\q}[1]{\begin{quote}\itshape #1\end{quote}}

\begin{document}

\begin{center}
{\LARGE\bfseries Lesson-plan-driven generation of\\personalised study packs}\\[6pt]
{\large Project status, and a request for raters}\\[4pt]
{\normalsize 5 August 2026}
\end{center}
\vspace{4pt}

\section{What the project does}

Most systems that generate teaching material take a curriculum objective, or a short topic
prompt, as the instruction. This project tests a richer instruction: the teacher's own
confirmed lesson plan. The plan becomes a binding specification. Everything generated after
it must stay inside that specification, inside the curriculum, and inside what the teacher
approves.

The prototype generates one kind of document: a take-home revision pack for Nigerian
Junior Secondary School~1 mathematics. The study uses five real lessons from one teacher,
with her lesson plans, her handwritten notes and her marked classwork.

\q{Does using a teacher-confirmed lesson plan, instead of curriculum objectives alone,
improve the quality of generated revision packs, and can those packs be personalised for
individual learners while staying within curriculum and teacher constraints?}

\subsection{The two hypotheses, and the conditions that test them}

\textbf{H1} — packs built from the lesson plan are rated higher than packs built from the
curriculum objectives alone. Grounding the pack in the teacher's source material, and
adding automated checking, improve it further.

\begin{center}
\begin{tabular}{@{}ll@{}}
\toprule
\textbf{C0} & Objectives only. No lesson plan, no source material.\\
\textbf{C1} & Add the teacher's lesson plan.\\
\textbf{C2} & Add the teacher's note as source material the pack must use.\\
\textbf{C3} & Add automatic checking and correction.\\
\textbf{C4} & Add teacher approval. Reported separately as governance evidence.\\
\bottomrule
\end{tabular}
\end{center}

\noindent The central comparison for H1 is \textbf{C1 against C0}.

\medskip
\textbf{H2} — a pack adapted to one learner's own evidence suits that learner better than
the approved pack the whole class receives, while still covering the lesson objectives.

\begin{center}
\begin{tabular}{@{}ll@{}}
\toprule
\textbf{P0} & The approved class pack, given to the learner unchanged.\\
\textbf{P1} & A pack adapted to that learner, from their marked classwork.\\
\bottomrule
\end{tabular}
\end{center}

\section{Where the project stands}

\begin{itemize}
\item \textbf{The system is built and runs end to end.} A confirmed lesson plan drives every
condition. Each run stores a reproducible record of what was generated and from what.

\item \textbf{Grounding works, and the effect is large.} When the teacher's note is supplied
as source material, the share of her material that reaches the pack rises from
$0.46$ to $0.83$. This is the clearest result so far.

\item \textbf{An automatic check is in place.} It needs no human and no language model. It
measures how much of the teacher's material reaches the pack, so a weak run is found before
expert time is spent on it.

\item \textbf{The lesson plan condition is not yet doing its job.} C1 does not yet beat C0.
This is the open problem, and Section~5 explains it.

\item \textbf{No expert ratings exist yet, by design.} The analysis plan was fixed on
26~July 2026, before any rating. The conditions can still be adjusted now. Once rating
starts they cannot.
\end{itemize}

\section{Timeline}

\begin{center}
\begin{tabular}{@{}p{32mm}p{100mm}@{}}
\toprule
\textbf{5--9 August} & Settle the C1 definition. Freeze all conditions. Re-run the
comparison and export the rating packets for H1.\\
\addlinespace
\textbf{9--10 August} & Run P0 and P1 for H2, across three lessons and their students.\\
\addlinespace
\textbf{From 11 August} & Expert rating begins. This depends on raters being recruited, and
on ethics approval.\\
\addlinespace
\textbf{10--21 August} & Two weeks' leave from work, taken to concentrate on writing the
dissertation. Rating runs in parallel: raters work through the queue independently.\\
\addlinespace
\textbf{From 22 August} & Analyse the returned ratings. Continue writing.\\
\bottomrule
\end{tabular}
\end{center}

\noindent What remains on my side is testing and final data checking, to confirm the
conditions are set correctly before they are frozen. The system itself is built.

\section{The request: expert raters}

Expert human ratings are the one thing I cannot complete myself.

\begin{center}
\begin{tabular}{@{}lrr@{}}
\toprule
& \textbf{Packs to rate} & \textbf{Ratings at 3 each}\\
\midrule
H1 (C0--C3) & 36 & 108\\
H2 (P0/P1) & 48--54 & 144--162\\
\midrule
\textbf{Total} & \textbf{84--90} & \textbf{about 250--270}\\
\bottomrule
\end{tabular}
\end{center}

\subsection{Do H1 and H2 need different raters? No — one pool serves both}

Both need the same kind of rater: an expert mathematics teacher, or a mathematics education
specialist. Each pack needs three independent ratings. \textbf{Five or six raters cover the
whole study}, because the work is distributed.

\begin{itemize}
\item \textbf{H1 is assigned by lesson, never by condition.} A rater who takes a lesson
rates all four versions of it. Splitting the work by condition would mix up how strict a
rater is with the condition being tested.
\item \textbf{H2 is assigned by learner.} Whoever rates P0 for a learner also rates P1, so
the comparison stays paired.
\item A shared set of four to eight packs, rated by everyone, measures how far the raters
agree with each other.
\end{itemize}

\subsection{The ask}

Could post-graduate education students, ideally in mathematics education, be recruited and
paid to rate? Alternatively, do you know practising or trainee mathematics teachers who
could take a block of packs? Raters score on 1--5 scales, and the packs hide which
condition produced them.

\clearpage

\section{The open problem: what the lesson plan condition should contain}

C1 is meant to be ``the lesson plan as the specification''. The difficulty is that this
teacher's lesson plan document contains very little teaching content. Almost everything she
teaches is in her separate handwritten note, which is what C2 supplies.

\subsection{What one lesson plan actually holds}

Lesson t3-w06, ``Construction'', is typical. The plan gives the subject, grade, term and
week; the topic and subtopic; four objectives; the entry behaviour; the materials
(a mathematical set); the textbook reference; and five steps. Each step gives a title, a
teaching pattern, a duration, and a sentence of narration.

That narration is the problem. It describes the classroom, not the mathematics:

\q{``Teacher explains the construction of line segments to the students while the students
construct line segments.''}

\noindent The pack is written for a learner revising alone at home, so it must not describe
the lesson as an event. Text of this kind was being supplied to the model and then ignored,
so it has been removed.

\subsection{What each condition sends}

\begin{center}
\begin{tabular}{@{}lcccc@{}}
\toprule
& \textbf{C0} & \textbf{C1} & \textbf{C2} & \textbf{C3}\\
\midrule
Subject, grade, topic, objectives & yes & yes & yes & yes\\
Step titles and their order & --- & yes & yes & yes\\
Teaching pattern of each step & --- & yes & yes & yes\\
Materials & --- & yes & yes & yes\\
The teacher's evaluation questions & --- & yes & yes & yes\\
The teacher's note as source material & --- & --- & yes & yes\\
Automatic checking and correction & --- & --- & --- & yes\\
\midrule
Classroom narration & \multicolumn{4}{c}{\emph{removed from every condition}}\\
\bottomrule
\end{tabular}
\end{center}

\subsection{Why C1 does not yet beat C0}

Two reasons, and the second is the serious one.

\begin{enumerate}
\item \textbf{The plan repeats what C0 already has.} A step title such as ``parallel and
perpendicular lines'' restates the objective. C0 already receives the full objective list,
so C1 adds structure the model can largely infer.

\item \textbf{The one substantial addition is missing from two of the five lessons.} The
teacher's own evaluation questions are real mathematics, and they do reach C1 --- but only
in three lessons. Two lessons have no evaluation step at all, so C1 gives them nothing
beyond a title and a duration.
\end{enumerate}

\noindent H1 needs C1 to beat C0 in every one of the five lessons. A condition that is
substantial in three lessons and thin in two cannot do that. On the automatic check, C1
currently leads in none of the five.

\subsection{What I propose}

Give C1 the position of each section in the lesson: that it is the second of four, what came
before it and what follows. The plan genuinely specifies this, it is present in all five
lessons, and it is not teaching content, so it does not blur the line between C1 and C2. It
also gives C1 something a rater can see --- sections that build on one another, rather than
four separate mini-lessons.

\clearpage

\section{Other open issues}

\begin{description}[leftmargin=0pt,style=nextline]
\item[The automatic critic's counts are too noisy to compare conditions.] Fifteen runs of
the same input gave counts ranging from zero to seven. The spread is larger than the
average. The counts are still useful for reading \emph{what} went wrong, but they cannot be
used to compare one condition against another. The check described in Section~2 is used for
that instead, and it is about fifty times more stable.

\item[A generator fault made earlier results unusable.] The generator writes a pack one
section at a time by calling tools. A library upgrade broke how it reads the model's
replies, and it silently failed on five of every six attempts. Only the runs that happened
to succeed were recorded, so any comparison from before 4~August is unreliable. The fault is
fixed: the current generator completed 100 runs out of 100.

\item[One teaching pattern suppresses practice questions.] Some steps are marked
``explanation only''. That describes how the teacher delivered the lesson, but the pack is
applying it to itself and producing fewer practice questions. A learner revising alone needs
practice whether or not the teacher set any that day.

\item[Correcting the teacher's material scores badly.] The automatic check rewards a pack
for carrying the teacher's material faithfully. If the pack corrects a mistake of hers, it
scores lower, because the corrected line no longer matches. This is a known limit of the
check, not a fault to be fixed, and it is one reason the expert ratings matter.

\item[One lesson's diagrams have no captions.] Fourteen figures in one lesson carry no
label, so the pack cannot refer to them by name. The other lessons are fully captioned.
\end{description}

\clearpage

\section{What H2 will look like}

H2 compares two packs for the same learner.

\begin{center}
\begin{tabular}{@{}p{18mm}p{110mm}@{}}
\toprule
\textbf{P0} & The approved class pack, reused unchanged for that learner. Nothing is
generated.\\
\addlinespace
\textbf{P1} & A pack adapted to that learner. The adaptation uses evidence taken from the
learner's own marked classwork. The classwork is written work, not multiple choice, so a
misconception is read from what the learner actually wrote.\\
\bottomrule
\end{tabular}
\end{center}

\noindent Three lessons have marked classwork. With eight or nine students each, that gives
48--54 packs and 144--162 ratings. The comparison is paired at the level of the
\emph{learner}, which is why H2 is well powered while H1 rests on only five lessons.

H2 depends on the teacher approving the class pack first, because both P0 and P1 are built
from it. Raters must see both packs \emph{before} any teacher edit, otherwise H2 measures
the teacher's editing rather than the personalisation.

\section{What a rating packet looks like}

\subsection{H1 --- one pack, judged against one lesson}

A rater receives three things.

\begin{enumerate}
\item \textbf{A lesson reference sheet.} Subject, grade, topic and subtopic; the objectives;
a short lesson sequence; the teacher's emphasis; the expected scope of the pack; and what is
explicitly out of scope. The same sheet is used for all four versions of that lesson, so
every version is judged against one target. The full lesson plan is deliberately not shown,
which limits bias.
\item \textbf{The generated pack, anonymised.} The rater does not know which condition
produced it. The order is fixed in advance and balanced across raters.
\item \textbf{The rating form.}
\end{enumerate}

\begin{center}
\begin{tabular}{@{}p{52mm}p{78mm}@{}}
\toprule
\textbf{H1 --- four scales, 1 to 5} & \textbf{What the rater judges}\\
\midrule
Objective and curriculum alignment & Covers the intended objectives and the curriculum target\\
Taught-lesson alignment & Reflects the lesson's sequence, emphasis and expected depth\\
Scope and correctness & Stays inside the lesson boundary, and avoids errors\\
Clarity and usefulness & Understandable and useful to a learner revising alone\\
\bottomrule
\end{tabular}
\end{center}

\subsection{H2 --- two packs for one learner, side by side}

The rater sees both packs for one learner, section by section, without being told which is
the adapted one.

\begin{center}
\begin{tabular}{@{}p{52mm}p{78mm}@{}}
\toprule
\textbf{H2 --- five scales, 1 to 5} & \textbf{What the rater judges}\\
\midrule
Suitability for the learner & Fits what this learner has shown they can do\\
Scaffolding & Gives the right amount of support\\
Challenge & Neither too easy nor too hard for this learner\\
Misconception support & Addresses the learner's specific errors (or not applicable)\\
Objective preservation & Still covers the lesson's original objectives\\
\bottomrule
\end{tabular}
\end{center}

\end{document}
